% 前言
\chapter*{前言}
\addcontentsline{toc}{chapter}{前言}

\section*{为什么学这门课？}

欢迎翻开《数学物理方法》这份讲义。不少同学第一次看到这门课的名字，难免会觉得它有点``硬''：既有数学，又有物理，还讲``方法''，仿佛注定公式很多、推导很长、门槛很高。

但不妨先把书名放在一边，从一个具体的问题出发：\textbf{为什么吉他琴弦能发出美妙的声音？}
一根弦、两端固定、轻轻拨动，这个动作看起来再普通不过，却已经包含了这门课最关心的几个问题：为什么会有不同的音高？为什么会出现不同的振动模式？边界条件为什么会决定系统的行为？如果再把真实乐器中的阻尼和能量耗散考虑进去，声音又为什么会逐渐减弱？

当你认真追问这些问题时，其实就已经站在数学物理方法的门口了。
在普通物理里，你早已接触过波动方程、热传导方程、拉普拉斯方程等对象。只是多数时候，我们更关心``会不会用''，还来不及系统追问：这些方程从哪里来？为什么写成这个样子？遇到不同的边界条件和初始条件，又该怎样求解、怎样理解？

这门课要做的，正是把这些零散出现的工具和观念串成一条清晰的线索，让你逐步学会\textbf{从物理问题出发，建立数学模型，完成数学求解，再回到物理解释}。
\section*{这门课学什么？}

如果把物理专业的学习过程看成一条不断向前延伸的链条，那么数学物理方法往往处在一个十分关键的位置。它向前承接高等数学和普通物理，向后连接理论力学、电动力学、量子力学、热力学与统计物理等课程，因此常被视为一门典型的桥梁课程。

\begin{figure}[H]
\centering
\begin{tikzpicture}[
    course/.style={rectangle, draw, rounded corners=3pt, minimum width=1.9cm, minimum height=0.75cm, align=center, font=\small, line width=0.8pt},
    arrow/.style={->, >=stealth, line width=0.9pt}
]
    % 前置课程（左侧）
    \node[course, fill=blue!8] (calc) at (-5, 3) {高等数学};
    \node[course, fill=blue!8] (mech) at (-5, 2) {力学};
    \node[course, fill=blue!8] (thermo) at (-5, 1) {热学};
    \node[course, fill=blue!8] (optics) at (-5, 0) {光学};
    \node[course, fill=blue!8] (em) at (-5, -1) {电磁学};
    \node[font=\footnotesize\bfseries, blue!60!gray] at (-5, 3.8) {基础课程};
    
    % 本课程（中间）
    \node[course, fill=orange!20, minimum width=2.8cm, minimum height=1.2cm, font=\small\bfseries] (mathphy) at (0, 1) {数学物理方法};
    \node[font=\footnotesize\bfseries, orange!60!gray] at (0, 2.2) {桥梁课程};
    
    % 后置课程（右侧）
    \node[course, fill=green!10] (tmech) at (5, 2.5) {理论力学};
    \node[course, fill=green!10] (edyn) at (5, 1.5) {电动力学};
    \node[course, fill=green!10] (qm) at (5, 0.5) {量子力学};
    \node[course, fill=green!10] (stat) at (5, -0.5) {热力学与统计物理};
    \node[font=\footnotesize\bfseries, green!50!gray] at (5, 3.3) {相关课程};
    
    % 箭头：前置 → 本课程
    \draw[arrow, blue!45] (calc.east) -- (mathphy.west);
    \draw[arrow, blue!45] (mech.east) -- (mathphy.west);
    \draw[arrow, blue!45] (thermo.east) -- (mathphy.west);
    \draw[arrow, blue!45] (optics.east) -- (mathphy.west);
    \draw[arrow, blue!45] (em.east) -- (mathphy.west);
    
    % 箭头：本课程 → 后置
    \draw[arrow, green!55!gray] (mathphy.east) -- (tmech.west);
    \draw[arrow, green!55!gray] (mathphy.east) -- (edyn.west);
    \draw[arrow, green!55!gray] (mathphy.east) -- (qm.west);
    \draw[arrow, green!55!gray] (mathphy.east) -- (stat.west);
\end{tikzpicture}
\caption{数学物理方法在课程体系中的位置}
\label{fig:course-position}
\end{figure}

换句话说，这门课并不是额外附加的一门技术课，而是把你前面学过的知识重新整理、提炼，并进一步衔接到后续专业课程中的一个枢纽环节。

\textbf{这些基础课程}提供了必要的数学基础和物理直觉：
\begin{itemize}
    \item \textbf{高等数学}：微积分、级数、常微分方程——这是本课程的数学基础。
    \item \textbf{普通物理}（力学、热学、光学、电磁学）：提供了丰富的应用背景。振动、波动、热传导、电磁场等问题，都会在本课程中被更系统地处理。
\end{itemize}

\textbf{后面的许多专业课程}会频繁用到这里学到的方法：
\begin{itemize}
    \item \textbf{理论力学}：拉格朗日方程、哈密顿方程的求解需要微分方程和变分法。
    \item \textbf{电动力学}：麦克斯韦方程组的求解离不开数学物理方程。
    \item \textbf{量子力学}：薛定谔方程的求解用到分离变量法、特殊函数。
    \item \textbf{热力学与统计物理}：不少模型的建立和分析，也需要本课程中的微分方程、积分变换等工具。
\end{itemize}

本课程大致可以分为四个彼此联系的部分：

\begin{figure}[H]
\centering
\begin{tikzpicture}[
    card/.style={draw, rounded corners=10pt, line width=0.9pt, minimum height=1.45cm, align=center, inner sep=7pt},
    arrow/.style={->, >=stealth, line width=0.95pt}
]
    % 顶部两个工具
    \node[card, fill=blue!8, draw=blue!45!black, text width=2.9cm] (complex) at (-2.8, 2.45)
    {\textbf{\large 复变函数}\\[2pt]\small 解析函数、级数展开、留数};
    \node[card, fill=green!8, draw=green!45!black, text width=2.9cm] (transform) at (2.8, 2.45)
    {\textbf{\large 积分变换}\\[2pt]\small 傅里叶变换、拉普拉斯变换};

    % 中间核心
    \node[card, fill=orange!12, draw=orange!60!black, text width=4.7cm, minimum height=1.65cm] (pde) at (0, -0.05)
    {\textbf{\large 数理方程}\\[3pt]\small 三类方程的建立与求解};

    % 底部产物
    \node[card, fill=purple!8, draw=purple!45!black, text width=4.8cm, minimum height=1.55cm] (special) at (0, -2.55)
    {\textbf{\large 特殊函数}\\[2pt]\small 勒让德、贝塞尔};

    % 连接箭头
    \draw[arrow, blue!45] ([xshift=0.55cm]complex.south) -- ([xshift=-1.1cm]pde.north);
    \draw[arrow, green!55!black] ([xshift=-0.55cm]transform.south) -- ([xshift=1.1cm]pde.north);
    \draw[arrow, purple!55!black] (pde.south) -- (special.north);
\end{tikzpicture}
\caption{课程四部分的逻辑关系}
\label{fig:course-structure}
\end{figure}

\begin{enumerate}
    \item \textbf{复变函数}：它把函数的视野从实轴扩展到复平面。进入这一部分以后，你会逐渐理解解析函数、级数展开、复积分和留数方法，也会看到复变观点在二维势场问题中所提供的几何直观。
    
    \item \textbf{积分变换}：它并不是凭空把问题``变没''，而是换一种表示、换一个观察角度，把原来难以下手的微分方程问题化成更容易处理的形式。傅里叶变换和拉普拉斯变换，是后面会反复使用的重要工具。
    
    \item \textbf{数理方程}：这是全书的核心部分。你会学习如何从具体物理过程出发建立偏微分方程，并掌握分离变量法、行波法、积分变换法、格林函数法等常用求解方法。前面学到的复变函数和积分变换，会在这里真正``用起来''。
    
    \item \textbf{特殊函数}：在求解数理方程的过程中，一些特殊函数会自然出现。它们不是额外增加的记忆负担，而是具有特定对称性的物理问题所对应的数学结果。球对称问题常常引出勒让德函数，圆柱对称问题常常引出贝塞尔函数。
\end{enumerate}

\medskip

如果这门课学得比较扎实，学完后你至少应该具备下面几方面的能力：

\begin{itemize}
    \item 建立并求解经典的偏微分方程（热传导方程、波动方程、拉普拉斯方程等）
    \item 熟练处理常见的边界条件和初始条件
    \item 理解复变函数在二维势场与复积分计算中的基本作用，并掌握留数法计算一些复杂实积分
    \item 能在球对称和圆柱对称问题中运用特殊函数求解具体问题
    \item 为后续理论物理课程打下更扎实的数学基础
    \item 面对新的物理问题时，知道应该优先调用哪些数学工具
\end{itemize}

更重要的是，你会逐渐熟悉一种典型的理论物理思路：\textbf{物理问题 $\rightarrow$ 数学模型 $\rightarrow$ 数学求解 $\rightarrow$ 物理解释}。这不仅适用于本课程，也会贯穿你后面的专业学习。

\section*{怎么学这门课？}

先把话说在前面：这门课确实不算轻松。它的难，不只在于公式多、推导长，更在于你需要在同一个问题里同时照看物理背景、数学结构和结果的物理意义。初学时，最容易卡住的往往是下面几类地方：

\begin{itemize}
    \item \textbf{复变函数}：从实数到复数的思维转换需要时间适应。复变函数的``导数''定义看起来和实函数相似，但条件要苛刻得多。学习时要多画图，尽量在复平面上理解它的几何意义，并把实变函数的直观和复变函数的结论对照起来看。
    
    \item \textbf{积分变换}：傅里叶变换和拉普拉斯变换的公式多，容易记混，收敛域又常常被忽略。学习时不要只背公式，要抓住它的本质：它是在换一种表示，使微分运算转化为更便于处理的形式。核心公式最好自己推几遍，把它们之间的联系真正理顺。
    
    \item \textbf{分离变量法}：公式本身并不复杂，真正的难点在边界条件。什么样的边界条件允许分离，分离以后常数如何确定，往往决定整道题能不能继续做下去。练习时最好从齐次边界条件起步，并养成先画边界示意图的习惯。
    
    \item \textbf{特殊函数}：名称多、性质多，刚开始很容易``晕头转向''。这一部分不要靠死记硬背硬顶过去，而要从它们的来源和物理意义入手，先弄清楚球对称问题为什么会引出勒让德函数，圆柱对称问题为什么会引出贝塞尔函数。
    
    \item \textbf{推导过程长}：很多定理证明和例题计算步骤多，最怕一路往下抄，却不知道自己到底在做什么。读长推导时，不妨强迫自己在关键地方停下来，问一句``为什么这一步可以这样做''。把逻辑主线抓住以后，后面的计算就不容易跟丢。
\end{itemize}

不过也不必被这些难点吓住。这门课还有一个很鲜明的特点：前面的工具会在后面的内容里不断重现。只要把几个核心方法真正弄懂，后面的学习反而会越来越顺。

\begin{itemize}
    \item \textbf{先问物理问题是什么}：每个数学概念都不是凭空出现的。先弄清楚它究竟想解决什么问题，再去看公式和定理，理解会稳很多。
    
    \item \textbf{把前修知识随时接上}：积分技巧、级数展开、常微分方程、向量分析这些内容都会反复用到。遇到卡点时，及时回头补，不要硬撑着往后赶。
    
    \item \textbf{坚持动手算}：很多步骤看上去已经明白了，真正下笔时问题才会暴露出来。数学物理方法终究是要靠自己一遍一遍算出来的，手感也是这样练出来的。
    
    \item \textbf{格外重视细节}：收敛域、边界条件、奇点位置、展开区间这些地方，往往决定一整道题能不能做对。
    
    \item \textbf{主动建立联系}：复变函数和二维势场、傅里叶变换和分离变量法、特殊函数和对称性，这些内容不是分开的。把联系看清楚，学习效率会高很多。
\end{itemize}

\begin{tishi}[动手算，才能真正掌握]
建议把每一道例题都当成一次``半开卷训练''：先独立推到自己真正卡住的地方，再回头对照讲义中的解答。不要只看最后结果，更要核对每一步为什么这样做；同时分清到底是概念不清、公式记混，还是边界条件处理失误。把问题找准，练习才会有效。
\end{tishi}


\section*{这份讲义怎么帮你学？}

课堂上，老师往往会不断提醒：哪里最容易出错，哪里要把前面的知识接回来，哪里又应该先抓住物理图像。这份讲义也希望承担这样的角色。

讲义中的\textcolor{red!70!black}{红色警示框}标注易错点，\textcolor{hintteal}{青色提示框}提供学习建议，\textcolor{physicsbrown}{棕色背景框}补充物理背景。它们不是装饰性的版式元素，而是帮助你抓住重点、放慢脚步、理清思路的阅读路标。

这份讲义从一开始就有一个非常明确的目标：\textbf{尽量写成一份真正适合学生自学的讲义}。它希望你在课堂上能跟着老师走，也希望你在课后独自翻开它时，依然知道这一页在讲什么、为什么这样讲，以及它和前后内容之间有什么关系。

编写时，编者始终默认读者的基础未必特别扎实，但愿意认真学，也愿意一步一步往前走。因此，这份讲义宁可把关键台阶讲得慢一点、细一点，也不愿意只把最后的结论摆在你面前。

具体来说，这份讲义在写法上会尽量做到下面几点：

\begin{itemize}
    \item \textbf{先讲来由，再讲结论}：抽象概念尽量先从物理问题、几何图像或典型例子引入，让你先知道``为什么会出现这个对象''，再去看正式定义和结论。
    
    \item \textbf{把关键台阶讲完整}：重要公式、定理和方法会尽量把推导过程写完整，把真正关键的环节解释清楚，让你知道每一步为什么成立，而不是只记住最后一个结果。
    
    \item \textbf{尽量用课堂语言讲清楚}：编者希望它读起来像老师在黑板前一步步讲，而不是像一本只给结论的手册。能直说的地方尽量不绕，能用直观语言说明的地方尽量不用过分生硬的表述。
    
    \item \textbf{在需要的地方补前修知识和图像}：遇到与高等数学、普通物理、常微分方程相关的内容，会补充必要的提醒、回顾和解释；复平面、积分路径、边界条件、振动模式等抽象对象，也会尽量配上示意图，帮助你把公式和图像对应起来。
    
    \item \textbf{帮助你把知识串成整体}：不仅介绍单个结论，也尽量提示不同章节之间的联系，让你看到一门课的整体脉络，而不是把它学成一堆零散技巧。
    
    \item \textbf{让复习和回顾都有抓手}：每章都安排了例题、小结和分层习题，方便你学完以后及时巩固，也方便考前回头整理思路。
\end{itemize}

如果你在阅读这份讲义时，感到自己不是在独自硬啃一堆零散结论，而是像有一位老师始终陪在身旁，把问题慢慢讲开，把难点层层拆开，把方法前后串起来，那么这份讲义就达到了它最初的目的。
\section*{致谢}

本讲义在编写过程中，参考了国内外多种优秀教材与资料，在此一并向有关作者致谢。

由于编者水平有限，书中难免仍有疏漏和错误之处，恳请读者批评指正。
\vspace{2em}
\begin{flushright}
\begin{tabular}{c}
编者 \\
2026年3月 \\
于湖南师范大学天文馆
\end{tabular}
\end{flushright}
